O material utilizado para a fabricação do eixo do tambor é o aço SAE 1045, cujas tensões admissíveis são:

\begin{table}[H]
  \centering
  \caption{Tensões admissíveis do eixo do tambor}
  \label{tab:tensoes_eixo_do_tambor}
  \begin{tabularx}{\textwidth}{>{\centering\arraybackslash}X>{\centering\arraybackslash}X>{\centering\arraybackslash}X}
    \toprule
    Descrição             & Símbolo    & Valor                  \\
    \midrule
    Tensão de escoamento  & $\sigma_e$ & \SI{400}{\mega\pascal} \\
    Tensão de resistência & $\sigma_R$ & \SI{600}{\mega\pascal} \\
    \bottomrule
  \end{tabularx}
\end{table}

A tensão de flexão admissível do eixo do tambor é dada pela equação \eqref{eq:tensao_admissivel_tambor}.

Substituindo os valores na equação \eqref{eq:tensao_admissivel_tambor}:
\begin{align*}
  \sigma_{\text{adm}} & = \frac{600}{1 \cdot \num{2.8}}           \\
  \sigma_{\text{adm}} & = \boxed{\SI{214.29}{\mega\pascal}}
\end{align*}

A tensão de cisalhamento admissível do eixo do tambor é dada pela equação \eqref{eq:tensao_cisalhamento_admissivel_tambor}.

Substituindo os valores na equação \eqref{eq:tensao_cisalhamento_admissivel_tambor}:

\begin{align*}
  \tau_{\text{adm}} & = \frac{600}{\sqrt{3}}                \\
  \tau_{\text{adm}} & = \boxed{\SI{123.72}{\mega\pascal}}
\end{align*}

\subsubsection{Pré-dimensionamento a flexo-torção}

A tensão de flexão do eixo do tambor é dada pela equação:

\begin{equation}
  \sigma_f = \frac{M_{f_e}}{W_f}
  \label{eq:tensao_flexao_eixo_do_tambor}
\end{equation}

Onde:

\begin{itemize}
  \item $\sigma_f$: tensão de flexão em \si{\mega\pascal}.
  \item $M_{f_e}$: momento fletor no eixo do tambor em \si{\newton\meter} (Tabela \ref{tab:dados_esforcos_tambor}).
  \item $W_f$: módulo de resistência à flexão em \si{\meter^3}.
\end{itemize}

\vspace{1em}

O módulo de resistência à flexão é dado pela equação:

\begin{equation}
  W_f = \frac{\pi d^3}{32}
  \label{eq:modulo_de_resistencia_a_flexao_eixo_do_tambor}
\end{equation}

Onde:

\begin{itemize}
  \item $W_f$: módulo de resistência à flexão em \si{\meter^3}.
  \item $d$: diâmetro do eixo do tambor em \si{\meter}.
\end{itemize}

\vspace{1em}

A tensão de cisalhamento por torção no eixo do tambor é dada pela equação:

\begin{equation}
  \tau_T = \frac{T_s}{W_T}
  \label{eq:tensao_cisalhamento_por_torcao_eixo_do_tambor}
\end{equation}

Onde:

\begin{itemize}
  \item $\tau_T$: tensão de cisalhamento por torção em \si{\mega\pascal}.
  \item $T_s$: torque de saída do redutor em \si{\newton\meter} \eqref{eq:torque_saida_redutor}.
  \item $W_T$: módulo de resistência à torção em \si{\meter^3}.
\end{itemize}

\vspace{1em}

O módulo de resistência à torção é dado pela equação:

\begin{equation}
  W_T = \frac{\pi d^3}{16}
  \label{eq:modulo_de_resistencia_a_torcao_eixo_do_tambor}
\end{equation}

Onde:

\begin{itemize}
  \item $W_T$: módulo de resistência à torção em \si{\meter^3}.
  \item $d$: diâmetro do eixo do tambor em \si{\meter}.
\end{itemize}

\vspace{1em}

Os módulos de resistência possuem a seguinte relação:

\begin{equation}
  W_T = 2 \cdot W_f
  \label{eq:relacao_modulos_eixo}
\end{equation}

A tensão combinada é dada pela equação:

\begin{equation}
  \sigma_c = \frac{M_c}{W_f} = \sqrt{{\sigma_f}^2 + 3 {\tau_T}^2}
  \label{eq:tensao_combinada_eixo_do_tambor}
\end{equation}

Onde:

\begin{itemize}
  \item $\sigma_c$: tensão combinada em \si{\mega\pascal}.
  \item $M_c$: momento fletor combinado em \si{\newton\meter}.
  \item $W_f$: módulo de resistência à flexão em \si{\meter^3}.
  \item $\sigma_f$: tensão de flexão em \si{\mega\pascal}.
  \item $\tau_T$: tensão de cisalhamento por torção em \si{\mega\pascal}.
\end{itemize}

\vspace{1em}

Desenvolvendo a equação \eqref{eq:tensao_combinada_eixo_do_tambor}, tem-se:

\begin{equation*}
  \frac{M_c}{\cancel{W_f}} = \sqrt{{\left(\frac{M_{f_e}}{\cancel{W_f}}\right)}^2 + 3 {\left(\frac{T_s}{2 \cancel{W_f}}\right)}^2}
\end{equation*}

\begin{equation}
  M_c = \sqrt{{M_{f_e}}^2 + \frac{3}{4} \cdot {T_s}^2}
  \label{eq:momento_fletor_combinado_eixo_do_tambor}
\end{equation}

Define-se que a tensão combinada deve ser menor ou igual à tensão de flexão admissível do material, ou seja:

\begin{equation*}
  \sigma_c \le \sigma_{\text{adm}}
\end{equation*}

Com isso, é possível obter a equação para calcular o diâmetro do eixo do tambor:

\begin{align*}
   & \sigma_{\text{adm}} = \frac{M_c}{W_f}                                        \\
   &                                                                              \\
   & \sigma_{\text{adm}} = \frac{M_c}{\frac{\pi d^3}{32}}                         \\
   &                                                                              \\
   & \sigma_{\text{adm}} = \frac{32 M_c}{\pi d^3}                                 \\
   &                                                                              \\
   & d^3 = \frac{32}{\pi} \cdot \frac{M_c}{\sigma_{\text{adm}}}                   \\
   &                                                                              \\
   & d = \sqrt[3]{\frac{32}{\pi}} \cdot \sqrt[3]{\frac{M_c}{\sigma_{\text{adm}}}}
\end{align*}

\begin{equation}
  d_{\text{eixo}} \approx \num{2.17} \cdot \sqrt[3]{\frac{M_c}{\sigma_{\text{adm}}}}
  \label{eq:diametro_eixo_do_tambor}
\end{equation}

\vspace{1em}

Substituindo os valores na equação \eqref{eq:momento_fletor_combinado_eixo_do_tambor}:

\begin{align*}
  M_c & = \sqrt{\num{2785.41}^2 + \frac{3}{4} \cdot \num{9150.44}^2} \\
  M_c & = \boxed{\SI{8399.79}{\newton\meter}}
\end{align*}

Substituindo os valores na equação \eqref{eq:diametro_eixo_do_tambor}:

\begin{align*}
  d_{\text{eixo}} & = \num{2.17} \cdot \sqrt[3]{\frac{\num{8399.79}}{\num{214.29} \cdot 10^6}} \\
  d_{\text{eixo}} & = \SI{0.07371}{\meter}                                        \\
  d_{\text{eixo}} & = \boxed{\SI{73.71}{\milli\meter}}
\end{align*}
